\documentclass[10pt,a4paper]{article}
\usepackage{amsmath,amssymb,amsthm,graphicx}
% Card-local minimal preamble (journal unit sigma-2025-003-comomentum-sections-and-poisson-maps-in-hami, P2-ahl16 precedent).
%% Verbatim from the pinned source: the author theorem environments and macro block,
%% so every macro and environment of the counted spans is available.
% Omitted and why: the documentclass line (the journal class sigma is not on the host;
%% the assembler supplies the class) and page-geometry/font blocks (they fight the
%% forced xelatex engine).
\usepackage{amsmath,amssymb,amsthm,mathtools}

\numberwithin{equation}{section}
\newtheorem{Theorem}{Theorem}[section]
\newtheorem*{Theorem*}{Theorem}
\newtheorem{Corollary}[Theorem]{Corollary}
\newtheorem{Lemma}[Theorem]{Lemma}
\newtheorem{Proposition}[Theorem]{Proposition}
 { \theoremstyle{definition}
\newtheorem{Definition}[Theorem]{Definition}
\newtheorem{Note}[Theorem]{Note}
\newtheorem{Example}[Theorem]{Example}
\newtheorem{Remark}[Theorem]{Remark}
\newtheorem{Question}[Theorem]{Question}
 }
\newcommand{\nom}{\nonumber}
\newcommand{\vph}{\boldsymbol{\varphi}}
\newcommand{\vsi}{\boldsymbol{\sigma}}
\newcommand{\bx}{\boldsymbol{x}}
\newcommand{\bq}{\boldsymbol{q}}
\newcommand{\bp}{\boldsymbol{p}}
\newcommand{\bC}{\boldsymbol{C}}
\newcommand{\bX}{\boldsymbol{X}}
\newcommand{\bJ}{\boldsymbol{J}}
\newcommand{\vrh}{\boldsymbol{\rho}}
\newcommand{\val}{{\boldsymbol{\alpha}}}
\newcommand{\vbe}{\boldsymbol{\beta}}
\newcommand{\vaa}{\boldsymbol{a}}
\newcommand{\vab}{\boldsymbol{b}}
\newcommand{\vom}{\boldsymbol{\omega}}
\newcommand{\veta}{\boldsymbol{\eta}}
\newcommand{\vxi}{\boldsymbol{\xi}}
\newcommand{\vchi}{{\boldsymbol{\chi}}}
\newcommand{\del}{\partial}
\newcommand{\ti}{\tilde}
\newcommand{\olv}{\overleftarrow}
\newcommand{\orv}{\overrightarrow}
\def\bF{{\bm{F}}}
\def\bG{{\bm{G}}}
\def\bg{{\bf g}}
\def\ba{{\bm{A}}}
\def\bb{{\bm{B}}}
\def\bu{{\bm{u}}}
\def\bv{{\bm{v}}}
\def\by{{\bm{Y}}}
\def\bz{{\bm{Z}}}
%\def\bx{\mbox{$x$}}
%\def\bxi{\mbox{$\xi$}}
\def\bxi{{\bm{\xi}}}
\def\bq{{\bm{q}}}
\def\bp{{\bm{p}}}
\def\bqk{\mbox{$q_k$}}
\def\bpk{\mbox{$p_k$}}
\def\bbx{{\bm{x}}}
\def\bbxi{{\bm{\xi}}}
\def\bbq{{\bm{q}}}
\def\bbp{{\bm{p}}}
\def\bbe{{\bm{e}}}
\def\bbc{{\bm{c}}}
\def\bbt{{\bm{t}}}
%\def\bbd{{\bm{d}}}
\def\bbF{{\bm{F}}}
\def\bbu{{\bm{u}}}
\def\bbv{{\bm{v}}}
\def\bbw{{\bm{w}}}
\def\bbz{{\bm{z}}}
\def\brho{{\bm{\rho}}}
\def\bzeta{{\bm{\zeta}}}
\newcommand{\vThe}{{\boldsymbol{\Theta}}}
\def\balpha{{\bm{\alpha}}}
\def\bomega{{\bm{\omega}}}
\def\bvartheta{{\bm{\vartheta}}}
\def\bOmega{{\bm{\Omega}}}
\def\bbone{{\bm{1}}}
\def\bTheta{{\bm{\Theta}}}
\def\bR{{\mathbb{R}}}
\def\bC{{\mathbb{C}}}
\def\bZ{{\mathbb{Z}}}
\def\bT{{\bm{T}}}
\def\bPhi{{\bm{\Phi}}}
\def\bh{{\bm{h}}}
\def\bphi{{\bm{\phi}}}
\def\bvarphi{{\bm{\varphi}}}
\def\bPsi{{\bm{\Psi}}}
\def\bdelta{{\bm{\delta}}}
\def\bchi{{\bm{\chi}}}
\newcommand{\calA}{{\mathcal A}}
\newcommand{\calB}{{\mathcal B}}
\newcommand{\calC}{{\mathcal C}}
\newcommand{\calD}{{\mathcal D}}
\newcommand{\calE}{{\mathcal E}}
\newcommand{\calF}{{\mathcal F}}
\newcommand{\calG}{{\mathcal G}}
\newcommand{\calH}{{\mathcal H}}
\newcommand{\calI}{{\mathcal I}}
\newcommand{\calJ}{{\mathcal J}}
\newcommand{\calK}{{\mathcal K}}
\newcommand{\calL}{{\mathcal L}}
\newcommand{\calM}{{\mathcal M}}
\newcommand{\calN}{{\mathcal N}}
\newcommand{\calO}{{\mathcal O}}
\newcommand{\calP}{{\mathcal P}}
\newcommand{\calR}{{\mathcal R}}
\newcommand{\calS}{{\mathcal S}}
\newcommand{\calT}{{\mathcal T}}
\newcommand{\calX}{{\mathcal X}}
\newcommand{\calY}{{\mathcal Y}}
\newcommand{\calU}{{\mathcal U}}
\newcommand{\calV}{{\mathcal V}}
\newcommand{\calW}{{\mathcal W}}
\newcommand{\Map}{{\rm Map}}
\newcommand{\ev}{{\rm ev}}
\newcommand{\diver}{{\rm div}}
\newcommand{\rdif}{\overleftarrow{\partial}}
\newcommand{\ld}{\overrightarrow{\partial}}
\newcommand{\sbv}[2]{{\{{{#1},{#2}}\}}}
\newcommand{\courant}[2]{{[{{#1},{#2}}]_D}}
\newcommand{\couranta}[2]{{[{{#1},{#2}}]_{DA}}}
\newcommand{\couranth}[2]{{[{{#1},{#2}}]_H}}
\newcommand{\courantr}[2]{{[{{#1},{#2}}]_D^{\pi}}}
\newcommand{\couranttp}[2]{{[{{#1},{#2}}]_D^{tP}}}
\newcommand{\liebra}[2]{{[{{#1},{#2}}]_{TM \oplus \bR}}}
\newcommand{\rhoe}{\rho_E}
\newcommand{\rhoa}{{\rho_A}}
\newcommand{\rhog}{{\rho_{\mathfrak{g}}}}
\newcommand{\rhott}{\rho_{T \oplus T^*}}
\newcommand{\rhota}{\rho_{T^*A^*}}
\newcommand{\rhotr}{\rho_{TM \oplus \bR}}
\newcommand{\rhoaa}{\rho_{A \oplus A^*}}
\newcommand{\bracket}[2]{\langle #1,\,#2\rangle}
\newcommand{\bracketa}[2]{{\langle #1,\,#2\rangle_A}}
\newcommand{\bracketm}[2]{{\langle #1,\,#2\rangle_{-}}}
\newcommand{\bbracket}[2]{\Biggl\langle #1,\,#2\Biggr\rangle}
\newcommand{\inner}[2]{{({{#1},{#2}})}}
\newcommand{\inwedbra}[2]{\langle #1\inwedge\,#2\rangle}
\newcommand{\cbral}{[\![}
\newcommand{\cbrar}{]\!]}
\newcommand{\vbracket}[2]{\cbral #1,\,#2\cbrar_D}
\newcommand{\excourant}[2]{\cbral #1,\,#2\cbrar}
\newcommand{\be}{\boldsymbol{e}}
\newcommand{\floor}[1]{{\lfloor #1 \rfloor}}
\newcommand{\variation}{{\hat{\delta}}}
\newcommand\Hom{\mathop{\mathrm{Hom}}\nolimits}
\newcommand\Conj{\mathop{\mathrm{Conj}}\nolimits}
\newcommand\Fun{\mathop{\mathrm{Fun}}\nolimits}
\def\brst{\mbox{$\delta$}}
\newcommand{\inwedge}{\stackrel\wedge,}
\newcommand{\rack}{\rhd}
\newcommand{\prack}{\trianglerighteq}
\newcommand{\cat}{\rightrightarrows}
\newcommand{\proj}{{pr}}
\newcommand{\pj}{{\pi}}
\newcommand{\dbrst}{{d_{BRST}}}
\newcommand{\md}{{d}}
\newcommand{\usigma}{\underline{\sigma}}
\newcommand{\ude}{{\underline{e}}{}}
\newcommand{\udx}{{\underline{\rd x}}{}}
\newcommand{\upartial}{{\underline{\partial}}{}}
\newcommand{\Bisec}{{\mathcal{B}(\calG)}}
\newcommand{\omegac}{{\omega_{can}}}
\newcommand{\omegam}{{\omega_{\calM}}}
\newcommand{\omegan}{{\omega_{\calN}}}
\newcommand{\Thetam}{{\Theta_{\calM}}}
\newcommand{\Thetan}{{\Theta_{\calN}}}
\newcommand{\rd}{\mathrm{d}}
\def\sd{{\bm{\mathrm{d}}}}
\def\sdd{{\bm{\mathrm{D}}}}
\def\bbd{{\bm{\mathrm{d}}}}
\newcommand{\lbracket}[2]{\cbral #1,\,#2\cbrar}
\def\xzero{X{}}
\newcommand{\rDD}{\mathrm{D}}
\newcommand{\uu}{{\underline{u}}{}}
\newcommand{\uv}{{\underline{v}}{}}
\newcommand{\ue}{{\underline{e}}{}}
\newcommand{\uf}{{\underline{f}}{}}
\newcommand{\us}{{\underline{s}}{}}
\newcommand{\ualpha}{{\underline{\alpha}}{}}
\newcommand{\ubeta}{{\underline{\beta}}{}}
\newcommand{\uJ}{{\underline{J}}{}}
\newcommand{\ua}{{\underline{a}}{}}
\newcommand{\ub}{{\underline{b}}{}}
\newcommand{\umu}{\underline{\mu}}
\newcommand{\unablamu}{\underline{\nabla \mu}}
\newcommand{\momega}{{\omega}}
\newcommand{\gomega}{{\omega}_{grad}}
\newcommand{\ggomega}{{\omega_{AKSZ}}}
\newcommand{\gS}{{S_{AKSZ}}}
\newcommand{\hmu}{{\widehat{\mu}}}
\newcommand{\tH}{{\tilde{H}}}
\newcommand{\rde}{{{}^E \rd}}
\newcommand{\rddr}{{\rd_{\mathrm{dR}}}}
\newcommand{\hrho}{{\widehat{\rho}}}
\newcommand{\hE}{{\widehat{E}}}
\newcommand{\baS}{{{}^A S}}
\newcommand{\rdd}{{\mathrm{D}}}
\newcommand{\rda}{{\mathrm{D}_A}}
\newcommand{\rdv}{{\mathrm{D}_V}}
\newcommand{\rdvn}{{\mathrm{D}_V^{\nabla}}}
\newcommand{\rdma}{{\mathrm{D}_{\mfa}}}
\newcommand{\mfa}{{\mathfrak{a}}}
\newcommand{\wpi}{{\widetilde{\pi}}}
\newcommand{\muvee}{{\mu^{\vee}}}
\newcommand{\bbT}{{\mathbb{T}}}
\newcommand{\nablabas}{{{}^A \nabla}^{\mathrm{bas}}}
\newcommand{\anabla}{{{}^A \nabla}}
\title{The basic curvature paired with a momentum section vanishes under the Poisson anchor}
\author{Hirota, Yuji, Ikeda, Noriaki}
\date{doi:10.3842/SIGMA.2025.003; CC BY 4.0}
\begin{document}
\maketitle
\setcounter{section}{3}
\setcounter{Theorem}{0}
\setcounter{equation}{2}
\noindent\textit{Editorial note.} All delivered mathematical text is the author's own, copied from the pinned journal source file sources/sigma-2025-003-comomentum-sections-and-poisson-maps-in-hami/pinned-original-source.tex (Comomentum Sections and Poisson Maps in Hamiltonian Lie Algebroids); no statement or proof step was rewritten, reordered or removed.\par\medskip
\noindent\textit{Editorial note.} The counted proof is the author's complete argument: All eleven equation labels invoked inside the span (LAidentity1, Poissonidentity, lcAtosion, lcbasiccurv, Pms12, Pms13, LAIden111, identity11, P3covlocal, identity12, identity13) are defined inside the span itself; the proof is self-contained and no external numbered reference is invoked (no cite command in the span). The Hamiltonian-Lie-algebroid axioms (P1)(P2)(P3) are restated in full inside the proof, so the external definition environment (source lines 739-743, which cites [Blohmann:2018]) is standing setup only; its citation lies outside the span and is disclosed here. The proof opens with 'They are proved by using local coordinate expressions' and carries two footnotes (Einstein summation convention; the TM-connection conventions) verbatim. The substitution and differentiation steps are delivered by result with the input identities named at each step (the source's terse computational style); carried verbatim. The card face renders Lemma 3.1 by source-text counter stepping (section 3, first shared-counter environment; seed section=3, Theorem=0, equation=2).\par\medskip
\noindent\textit{Editorial note.} The label basicms names the counted statement, and the proof environment immediately following it in the source is delivered verbatim as its proof.\par\medskip
\noindent\textit{Editorial note.} The counted statement is carried by the author's own Lemma environment in the article source.\par\medskip
\noindent\textit{Editorial note.} Journal-unit assembly: card-local preamble (author theorem environments and macros verbatim), article class, no bibliography in the counted span (the span carries no \textbackslash{}cite commands).\par\medskip
\noindent\textit{Editorial note.} Revision r002: regenerated from the same spec after the r001 assembly crashed during packaging.\par\medskip
\noindent\textit{Editorial scope.} the counted statement and its own proof environment only\par\medskip
\noindent\textit{Allowed standard prerequisites (editorial).}\begin{itemize}
\item A Lie algebroid (A, rho, [-,-]) over a smooth manifold M, with the local coordinate setup delivered inside the proof: coordinates x\textasciicircum{}i on M, fiber bases e\_a of A and dual e\textasciicircum{}a of A*, anchor expression rho(e\_a) = rho\_a\textasciicircum{}i(x) partial\_i, bracket [e\_a, e\_b] = C\_\{ab\}\textasciicircum{}c(x) e\_c, and the two Lie algebroid identities (LAidentity1) with the Cicly terms.
\item A Hamiltonian Lie algebroid over a Poisson manifold: the defining axioms (P1), (P2), (P3) are restated in full inside the proof (source lines 1033-1050) as equations (Pms12), (Pms13); the definition environment itself (source lines 739-743, label HamiltonianLAS) is standing setup and cites [Blohmann:2018] externally, outside the span.
\item A momentum section mu, a TM-connection nabla on A with connection 1-form omega, and the local coordinate expressions of the A-torsion, the curvature and the basic curvature (lcAtosion, lcbasiccurv), all delivered inside the proof; the Poisson condition [pi, pi]\_S = 0 (Poissonidentity) likewise in-span.
\end{itemize}\medskip
\begin{Lemma}\label{basicms}
In a Hamiltonian Lie algebroid over a Poisson manifold,
$\pi^{\sharp} \bigl\langle{}^A S,\mu\bigr\rangle =0$,
where $\bracket{-}{-} =0$ is the pairing of $A$ and $A^*$.
\end{Lemma}

\begin{proof}
They are proved by using local coordinate expressions.

For a Lie algebroid $(A, \rho, [-,-])$ over a smooth manifold $M$,
let $x^i$ be a local coordinate on~$M$, $e_a \in \Gamma(A)$ be a basis of the fiber of $A$ and $e^a \in \Gamma(A^*)$ be the dual basis on $A^*$.
The basis of~$TM$ is denoted by $\partial_i = \tfrac{\partial}{\partial x^i}$.
$i$, $j$, etc. are indices on $M$ and $a$, $b$, etc. are indices on
the fiber of~$A$ and~$A^*$.
Upper indices are ones for $TM$ and $A$, and lower indices are ones
for $T^*M$ and~$A^*$.\footnote{Einstein summation convention is used for
sums in local coordinate expressions.}

Local coordinate expressions of the anchor map and the Lie bracket are
$\rho(e_a)= \rho^i_a(x) \partial_i$ and~${[e_a, e_b ] = C_{ab}^c(x) e_c}$, where $\rho^i_a(x)$ and $C_{ab}^c(x)$
are local functions.
Identities of the Lie algebroid are\looseness=-1
\begin{gather}%\label{LAidentity2}
 \rho_a^j \partial_j \rho_{b}^i - \rho_b^j \partial_j \rho_{a}^i = C_{ab}^c \rho_c^i,
\qquad
 C_{ad}^e C_{bc}^d + \rho_a^i \partial_i C_{bc}^e + \operatorname{Cycl}(abc) = 0.\label{LAidentity1}
\end{gather}
The condition of Poisson bivector field $\pi$, $[\pi, \pi]_S=0$, is
\begin{gather}
 \pi^{il} \partial_l \pi^{jk} + \operatorname{Cycl}(ijk) = 0.
\label{Poissonidentity}
\end{gather}
A connection $1$-form $\omega_a^b = \omega_{ai}^b \rd x^i$
for the connection $\nabla$ is introduced as
$\nabla_i e_a := \omega_{ai}^b e_b$ for the basis, $e_a$.
Local coordinate expressions of an $A$-torsion,
a curvature and a basic curvature are given by
\begin{gather}%\label{lccurv}
T_{ab}^c =
- C_{ab}^c + \rho_a^i \omega_{bi}^c - \rho_b^i \omega_{ai}^c,
\label{lcAtosion}
\\
R_{ijb}^a =
\partial_i \omega_{aj}^b - \partial_j \omega_{ai}^b
+ \omega_{ai}^c \omega_{cj}^b - \omega_{aj}^c \omega_{ci}^b,
\nonumber
\\
S_{iab}^{c} =
\nabla_i T_{ab}^c + \rho_b^j R_{ija}^c - \rho_a^j R_{ijb}^c,
 \nonumber \\
\phantom{S_{iab}^{c} }{}= - \partial_i C^c_{ab} - \omega_{di}^c C_{ab}^d
+ \omega_{ai}^d C_{db}^c + \omega_{bi}^d C_{ad}^c
+ \rho_a^j \partial_j \omega_{bi}^c
- \rho_b^j \partial_j \omega_{ai}^c
\nonumber \\
\phantom{S_{iab}^{c} =}{}+ \partial_i \rho_a^j \omega_{bj}^c
- \partial_i \rho_b^j \omega_{aj}^c
+ \omega_{ai}^d \rho_d^j \omega_{bj} ^c
- \omega_{bi}^d \rho_d^j \omega_{aj} ^c,
\label{lcbasiccurv}
\end{gather}
where the covariant derivative $\nabla_i T_{ab}^c$ is given by
\[
\nabla_i T_{ab}^c =
\partial_i T_{ab}^c
+ \omega_{di}^c T_{ab}^d - \omega_{ai}^d T_{db}^c - \omega_{bi}^d T_{ad}^c.
\]

Equations~(P1), (P2) and (P3) in the definition of a Hamiltonian Lie algebroid are
\begin{gather}
{\rm (P1)} \quad {}^A \nabla_a \pi^{ij}
= \rho_a^k \partial_k \pi^{ij} - \partial_k \rho^i_a \pi^{kj}
+ \rho^i_b \omega^b_{ak} \pi^{kj}
- \partial_k \rho^j_a \pi^{ik}
+ \rho^j_b \omega^b_{ak} \pi^{ik} = 0,\nonumber\\
{\rm (P2)} \quad
\rho^i_a = \pi^{ij} \nabla_j \mu_a
\bigl(= \pi^{ij} \bigl(\partial_j \mu_a - \omega_{aj}^b \mu_b\bigr)\bigr),
\label{Pms12}\\
{\rm (P3)}\quad
\rho^i_{[a} \partial_i \mu_{b]} - C_{ab}^c \mu_c
= - \pi^{ij} \nabla_i \mu_a \nabla_j \mu_b
\bigl(= \pi^{ij} \bigl(\partial_i \mu_a - \omega_{ai}^c \mu_c\bigr)
\bigl(\partial_j \mu_b - \omega_{bj}^d \mu_d\bigr)\bigr).
\label{Pms13}
\end{gather}

Substituting \eqref{Pms12} into the identity \eqref{LAidentity1},
$\rho_a^j \partial_j \rho_{b}^i - \rho_b^j \partial_j \rho_{a}^i = C_{ab}^c \rho_c^i$,
we obtain the following equation:
\begin{gather}
 \pi^{ij} \bigl(-\pi^{kl} \partial_k \nabla_j \mu_a \nabla_l \mu_b
- \pi^{kl} \nabla_k \mu_a \partial_l \nabla_j \mu_b
- \partial_j \pi^{kl} \nabla_k \mu_a \nabla_l \mu_b
- C_{ab}^c \nabla_j \mu_c\bigr)=0.
\label{LAIden111}
\end{gather}
Here, we use $\pi$ is the Poisson bivector field, i.e.,
equation~\eqref{Poissonidentity}.
Moreover, equation~\eqref{LAIden111} is equivalent to
\begin{gather}
\pi^{ij} \bigl(\pi^{kl} \nabla_k \nabla_j \mu_a \nabla_l \mu_b
+ \pi^{kl} \nabla_k \mu_a \nabla_l \nabla_j \mu_b
+ \partial_j \pi^{kl} \nabla_k \mu_a \nabla_l \mu_b
- T_{ab}^c \nabla_j \mu_c\bigr)=0,
\label{identity11}
\end{gather}
which is the covariant expression of equation~\eqref{LAIden111}.

We consider another identity as follows.
Substituting \eqref{Pms12} to \eqref{Pms13}
and using equation~\eqref{Poissonidentity},
we obtain
\[%\label{mumuidentity}
 \pi^{ij} \nabla_j \mu_a \partial_i \mu_b
- \pi^{ij} \nabla_j \mu_b \partial_i \mu_a
+ \pi^{ij} \nabla_i \mu_a \nabla_j \mu_b
- C_{ab}^c \mu_c=0.
\]
The covariant expression of this identity is
\begin{gather}
\pi^{kl} \nabla_k \mu_a \nabla_l \mu_b
- T_{ab}^c \mu_c=0,
\label{P3covlocal}
\end{gather}
where equation~\eqref{lcAtosion} is used.
The global form of equation~\eqref{P3covlocal} is
\[
\pi(\nabla \mu(e_1), \nabla \mu(e_2))
- \big\langle\mu,{}^A T(e_1, e_2)\big\rangle = 0
\]
for $e_1, e_2 \in \Gamma(A)$.

Acting the covariant derivative by $\nabla_j$ to equation~\eqref{P3covlocal},
and using the identity \eqref{lcbasiccurv},
\smash{$S_{iab}^{c} = \nabla_i T_{ab}^c + \rho_b^j R_{ija}^c - \rho_a^j R_{ijb}^c$},
we obtain the equation
\begin{gather}
 \pi^{kl} \nabla_k \nabla_j \mu_a \nabla_l \mu_b
+ \pi^{kl} \nabla_k \mu_a \nabla_l \nabla_j \mu_b
+ \nabla_j \pi^{kl} \nabla_k \mu_a \nabla_l \mu_b
- T_{ab}^c \nabla_j \mu_c
- S_{jab}^c \mu_c =0.
\label{identity12}
\end{gather}

By taking $\pi^{ij} \times$ \eqref{identity12},
we obtain the identity
\begin{gather}
 \pi^{ij}\bigl(\pi^{kl} \nabla_k \nabla_j \mu_a \nabla_l \mu_b
+ \pi^{kl} \nabla_k \mu_a \nabla_l \nabla_j \mu_b
+ \nabla_j \pi^{kl} \nabla_k \mu_a \nabla_l \mu_b\nonumber\\
\qquad
{}- T_{ab}^c \nabla_j \mu_c
- S_{jab}^c \mu_c\bigr) =0.
\label{identity13}
\end{gather}
Compare equations~\eqref{identity11} and~\eqref{identity13},
where note that $\nabla_j \pi^{kl} = \partial_j \pi^{kl}$
since we do not introduce a $TM$-connection on $TM$.\footnote{We have introduced a $TM$-connection on $A$, $\nabla$ and
an $A$-connection on $TM$, \smash{$\anabla = \nablabas$}.}
Consistency between equations~\eqref{identity11} and~\eqref{identity13}
gives the following identity,
$\pi^{ij} S_{jab}^c \mu_c =0$,
i.e.,
$ \pi^{\sharp} \bracket{{}^AS}{\mu} =0$.
\end{proof}

\end{document}
